\documentclass[fontsize=10pt]{article}
\usepackage[margin=0.70in]{geometry}
\usepackage{lipsum,mwe,abstract}
\usepackage[T1]{fontenc}
\usepackage[english]{babel}

\usepackage{fancyhdr}
\setlength\parindent{0pt}

\usepackage{amsmath,amsfonts,amsthm}
\usepackage{wrapfig}
\usepackage{graphicx}
\usepackage{float}
\usepackage{subcaption}
\usepackage{comment}
\usepackage{enumitem}
\usepackage{cuted}
\usepackage{sectsty}
\usepackage{eurosym}
\usepackage{xcolor}
\usepackage{hyperref}
\hypersetup{colorlinks=true, urlcolor=blue!50!black, linkcolor=black}
\allsectionsfont{\normalfont \normalsize \scshape}

\newcommand{\figplaceholder}[2]{%
  \par\medskip
  \fcolorbox{black!50}{black!4}{%
    \begin{minipage}{\dimexpr#1-2\fboxsep-2\fboxrule\relax}
      \vspace{0.6em}
      \centering
      \textbf{\textcolor{red!70!black}{Figure source requested from S.~Pazos / M.~Lanza}} \\[0.4em]
      \footnotesize
      \emph{We would prefer the original measurement / vector / slide
      file for this figure; the description below is detailed enough
      to redraw if the source is not readily available.} \\[0.6em]
      \footnotesize\textit{Content:} #2
      \vspace{0.6em}
    \end{minipage}%
  }%
  \par\medskip
}

\renewenvironment{abstract}
 {\small
  \begin{center}
  \bfseries \abstractname\vspace{-.5em}\vspace{0pt}
  \end{center}
  \list{}{%
    \setlength{\leftmargin}{0mm}
    \setlength{\rightmargin}{\leftmargin}%
  }
  \item\relax}
 {\endlist}

\makeatletter
\renewcommand{\maketitle}{\bgroup\setlength{\parindent}{0pt}
\begin{flushleft}
  \textbf{\@title}
  \@author \\
  \@date
\end{flushleft}\egroup
}
\makeatother

%% -------------------------------------------------------------------

\title{
\Large NS-RAM Co-Design: A Differentiable Software Bridge from Silicon to Algorithm
}
\date{\today}
\author{E. Bergvall, R. Luciani, S. Pazos, M. Lanza}

\begin{document}

\maketitle

% --------------- ABSTRACT
\begin{abstract}
The NS-RAM 2T floating-body cell delivers synapse, neuron and
short-term-memory behaviour in unmodified $130$~nm bulk CMOS, with
published per-cycle energy in the sub-pJ range and a per-cell
footprint of tens of square micrometres. What is missing today is
the bridge between silicon characterisation and algorithm design:
each cell-parameter or routing-topology question round-trips through
hand-curated SPICE decks and fab cycles. We are building that
bridge --- a differentiable, GPU-batched Python port of the BSIM4
channel and a Gummel--Poon / VBIC parasitic NPN, wired into a
body-state-aware network simulator that scales from a single device
to hundreds of thousands of cells (Fig.~\ref{fig:codesign}).

The artifact already runs as a functional research tool. On the
full $25$-bias DC anchor set it reaches a cell-wide
$\log_{10}|I_d|$ RMSE of approximately $1.2$~decade
(forward$+$backward averaged, pseudo-transient solver),
demonstrably good enough that downstream network demonstrations
inherit silicon-realistic behaviour. Six of nine pre-registered
dynamic-validation tests pass on the calibrated cell --- DC fit,
forward/backward hysteresis, nanosecond drain-snap, body-voltage
latching, sub-threshold integration, and a clean firing threshold ---
while two transient targets (autonomous self-reset and free-running
relaxation oscillation) remain open and are the subject of the
near-term physics-completion plan. On top of the cell, seven
network applications pass all pre-registered evaluation gates,
spanning reservoir computing, hyperdimensional encoding, associative
binding, predictive coding, stochastic computing, adaptive
equalisation, and an edge sensing cascade.

This proposal funds the work that turns the functional artifact into
a silicon-calibrated one: close the transient gap with the
remaining body-circuit physics, extend the network demonstrations to
realistic device variation and shared-rail layout, and convert two
pending characterisation runs on existing silicon into a
tape-out-ready cell-parameter recommendation.
\end{abstract}

\begin{figure}[h]
  \centering
  \includegraphics[width=0.95\textwidth]{figures/codesign_loop.pdf}
  \caption{The co-design loop this proposal funds. Measured
  silicon (left) is the substrate; the differentiable Python
  simulator (centre) inverts target algorithmic primitives back into
  cell-parameter and routing-topology specifications; the next
  tape-out (right) consumes those specs. The dashed return path
  closes the loop: each fabricated batch re-feeds the simulator.
  Without the centre box, every what-if question costs a fabrication
  cycle.}
  \label{fig:codesign}
\end{figure}

\rule{\linewidth}{0.5pt}

% --------------- MAIN CONTENT

\section{Why NS-RAM, and why a differentiable simulator}

Unlike RRAM, PCM and memristor approaches, NS-RAM is fabricated in
\emph{standard bulk CMOS} using only intrinsic-silicon parasitics:
no exotic materials, no novel mask steps, no separate process
integration. The same combination of impact-ionisation, floating-body
charging and a parasitic lateral NPN that conventional designs treat
as a latch-up hazard is, in the NS-RAM topology, the operating
mechanism. Reported behaviours of the cell include sub-pJ
write/spike energies, nanosecond-scale drain-current spikes
($\sim$\,$4$--$5$~mA peak, tens-of-ns rise), driven oscillation in
the few-MHz range, and analog short-term memory that persists for
microseconds to milliseconds depending on body resistance. These
properties --- in plain CMOS --- are what make the cell a credible
substrate for energy-constrained neuromorphic inference.

The bottleneck is not silicon, it is iteration speed at the
algorithm level. A useful neuromorphic substrate is one a designer
can \emph{simulate, vary, and back-propagate through}. A
differentiable port lets reservoir spectral radii, hyperdimensional
encoding biases, and predictive-coding gains be specified as
algorithmic targets and inverted to cell parameters. It also lets
device variation be drawn from realistic distributions and folded
into network-level robustness studies before any mask is ordered.
That loop is what this proposal funds.

% ------------------
\begin{figure}[h]
  \centering
  \includegraphics[width=0.99\textwidth]{figures/brief_headlines_honest/brief_headlines.pdf}
  \caption{Where the artifact stands today. Left: DC anchor fit on
  the $25$-bias measurement set, cell-wide $\log_{10}|I_d|$ RMSE
  $\approx 1.2$~decade with the pseudo-transient solver, on full
  forward$+$backward sweeps. Centre: nine-test dynamic-behaviour
  battery; six pass (DC, hysteresis, nanosecond snap, latch,
  integrate, threshold), three open (self-reset, free oscillation,
  knee-gate). Right: seven network demonstrations cleared
  pre-registered gates on independent benchmarks.}
  \label{fig:headlines}
\end{figure}

\section{What we have built}

The artifact is a single Python package, GPU-batched through PyTorch
on commodity integrated graphics, that exposes three layers:

\begin{enumerate}[leftmargin=1.3em, topsep=2pt, itemsep=2pt]
\item \textbf{A differentiable BSIM4 channel.}
A port of the BSIM4 v4.8.3 charge and current equations against the
public PTM~130 model card, validated bias-by-bias against ngspice.
The port reproduces ngspice currents to within a fraction of a
decade across the operating box that matters for the cell, exposes
all gradients through PyTorch autograd, and runs at roughly
$5\!\times\!10^6$ cell-evaluations per second on an integrated AMD
Radeon~8060S GPU.

\item \textbf{A parasitic-NPN coupling.}
A Gummel--Poon NPN tied to the BSIM4 channel through impact
ionisation and floating-body charging, with a VBIC level-4 path
available as a swap-in. A pseudo-transient body integration is the
default DC solver and is what produces the best honest cell-wide
RMSE.

\item \textbf{A network simulator.}
A scaling backend that batches cell evaluations across topologies
(Erd\H{o}s--R\'enyi, small-world, ring, lattice, and reservoir
families), reaches $N\!\geq\!10^5$ cells on commodity hardware, and
exposes a standard graph-construction interface to downstream
applications. Cell variation, weight initialisation, and readout
training are all standard PyTorch operations.
\end{enumerate}

\section{Validation status, honest}

\subsection{DC anchor fit}
On the full set of $25$ measured biases (forward and backward
sweeps, no direction-dropping, no bias-dropping), the
pseudo-transient pipeline reaches a cell-wide
$\log_{10}|I_d|$~RMSE of approximately $1.19$~decade
(forward~$1.35$, backward~$1.03$). With the VBIC NPN this is
$\approx 1.28$~decade. These are the publishable headline numbers;
earlier ``$<1.0$~decade'' values in our private logs are
forward-only or cherry-bias subsets and we no longer report them.

\subsection{Dynamic-behaviour battery}
We pre-registered nine behaviours the cell must exhibit to count as
a faithful surrogate for published NS-RAM phenomenology.
\textbf{Six pass} on the calibrated cell:
DC characteristic at the publishable RMSE,
forward/backward hysteresis,
nanosecond drain-current snap on a step input,
floating-body latch into the elevated-$V_b$ regime,
sub-threshold integration of a sub-threshold pulse train,
and a sharp firing threshold.
\textbf{Three remain open:}
the knee-gate $V_{G2}$ regime selector (numerical issue, in repair),
the autonomous self-reset back to the resting body voltage
(physics: requires the impact-ionisation/body-charge feedback to
close on its own time-constant rather than be terminated by the
input), and a free-running relaxation oscillation in the
$\sim\!430$~ns target band that matches the published driven
spike-train waveforms.

\subsection{Network demonstrations}
Seven network-level applications were dispatched against pre-
registered gates on independent benchmarks (no benchmark selected
post-hoc to confirm the cell). \textbf{All seven pass:}
(i) a sparse Erd\H{o}s--R\'enyi reservoir on Mackey--Glass
($\tau\!=\!17$) reaching normalised RMSE~$0.0153$ at $N\!=\!1024$;
(ii) hyperdimensional encoding on UCI~HAR reaching $\approx 84\%$
test accuracy at $D\!=\!8192$;
(iii) an associative-memory binder at $\approx 87\%$ recall and
$32$-pair capacity;
(iv) predictive coding on the NAB anomaly-detection benchmark at
$F1\!=\!0.34$;
(v) a $1/f$-noise stochastic-computing RNG passing all five NIST
SP800-22 batteries we ran;
(vi) a $16$-tap NS-RAM LMS equaliser on a $3$-echo QPSK channel,
BER~$=\!0$ at $20$~dB and $\approx 1.6\!\times\!10^{-2}$ at $10$~dB,
at $\sim\!170\times$ lower energy per symbol than fp32 LMS;
(vii) a KWS$\to$ECG edge cascade reaching $F1\!=\!0.845$ at
$\approx\!61\%$ energy savings against an always-on baseline.
Each ran end-to-end on the same calibrated cell, with no per-task
re-tuning of the device parameters.

% ------------------
\begin{figure}[h]
  \centering
  \includegraphics[width=0.99\textwidth]{figures/network_in_action/network_in_action.pdf}
  \caption{Seven network demonstrations on the same calibrated
  cell, one panel each. Reservoir computing (Mackey--Glass),
  hyperdimensional encoding (UCI~HAR), associative-memory binding,
  predictive coding (NAB), stochastic computing ($1/f$ RNG),
  adaptive equalisation, and an edge sensing cascade. None required
  per-task tuning of the underlying device parameters; the variation
  is in graph topology, readout, and inputs.}
  \label{fig:netinaction}
\end{figure}

\section{What we are scoping next}

\subsection{Close the transient gap}
The two open dynamic tests --- self-reset and relaxation
oscillation --- share a root cause: the body-discharge time-constant
in the lumped model is set by junction capacitance alone, which is
$\sim\!100\times$ too small to match published recovery times for
the cell.  Recent literature on parasitic-NPN dynamics in DNW CMOS
and on NS-RAM-class operating regimes points to two physical
additions: (i)~a distributed body-resistance network (the
$10$~k$\Omega$--$1$~M$\Omega$ tunable $R_B$ band documented for the
cell), and (ii)~a body-capacitance contribution from interface and
inversion-layer charge that is not captured by the junction-only
$C_j$ in the model card. Both can be added inside the existing
differentiable port; the work item is parameter identification, not
new infrastructure. Closing the transient gap promotes the artifact
from a DC-faithful network simulator to a transient-faithful one.

\subsection{Network demonstrations under realistic noise}
The seven passing demonstrations use the calibrated cell at nominal
parameters. The next step is to fold realistic device-to-device
variation, shared-bulk-rail coupling, and silicon-grade $1/f$ and
shot noise into the network simulator and re-evaluate.  This is the
study that tells a designer whether a particular topology
recommendation survives fabrication.

\subsection{Path to silicon validation}
Two existing-silicon characterisation runs convert the present
fit-defensible parameter point into a silicon-calibrated one: an
$I_c/I_b$ saturation extraction for the lateral parasitic NPN, and
a pulsed-$V_d$ or transmission-line-pulse measurement of the
body-region $R_b\!\cdot\!C_b$ time constant. With those two
measurements in hand, the body-state surrogate is pinned to silicon
and the cell-parameter recommendations the simulator produces become
tape-out-actionable.

\section{Deliverables and timeline}

\textbf{M3} (achieved): differentiable cell, network simulator,
$\approx 1.2$~decade cell-wide DC fit.
\textbf{M6} (Sep~2026): transient gap closed (nine-of-nine
dynamic-behaviour battery passing), network demonstrations
re-evaluated under realistic noise, full benchmark suite at
$N\!\in\!\{10,10^2,10^3,10^4\}$.
\textbf{M9} (Dec~2026): fan-out and shared-bulk-rail validation
circuit ($10$--$30$ cells), feeding a tape-out-ready cell-parameter
recommendation. \textbf{M12} (Mar~2027): thick-oxide neuron-cell
variant modelled and benchmarked, completing the
synapse--soma--linear-input cell family.

\section{Budget}
Total request approximately \textbf{$200$~k\euro} for $12$ months
(extension optional, conditional on the M6 review).
Personnel: Bergvall $0.8$~FTE ($140$~k\euro);
Luciani $0.3$~FTE optional ($35$~k\euro, Julia/Enzyme cross-
validation collaborator); Pazos and Lanza in-kind via KAUST.
Travel $9$~k\euro\ (one KAUST tape-out review trip plus one
neuromorphic-hardware venue); $5$~k\euro\ hardware contingency;
$3$~k\euro\ open-access and replication; institutional indirect
$\approx 5\%$. Compute is already in place
(HP~Z2 + AMD~Ryzen~AI~Max+~$395$ with Radeon~$8060$S integrated
GPU; $5\!\times\!10^6$ cell-evaluations per second confirmed; the
GPU port handles reservoirs up to $N\!=\!2\!\times\!10^4$ via a
block-matmul that bypasses a vendor driver bug at larger scales).

\vspace{2pt}
\hrule
\vspace{2pt}

{\footnotesize\emph{What we will defend in front of any reviewer.}
(i)~a differentiable, GPU-batched Python port of the cell, validated
against ngspice; (ii)~a cell-wide DC fit of approximately
$1.2$~decade on the full bias set, forward$+$backward, no
direction-pick; (iii)~six of nine dynamic behaviours passing on the
calibrated cell, with the two open transient targets explicitly
named; (iv)~seven network demonstrations passing pre-registered
gates on independent benchmarks; (v)~a concrete, two-measurement
silicon-validation plan that turns the artifact into a tape-out
recommendation.
\emph{What we will not defend.} ``NS-RAM is a better reservoir than
software ESN'' (false on like-for-like benchmarks --- our claim is
the energy-per-inference floor, not algorithmic superiority); any
$<1.0$~decade DC headline (those are forward-only or cherry-bias
subsets); a fully autonomous relaxation oscillator on the present
cell (that target is open and is the first item on the M6 list).}

\end{document}
